\chapter{Минимальная typed-архитектура бинарного \texorpdfstring{$u/d$}{u/d}-различителя}
\label{ch:v8-up-down-minimal-typed-discriminator}

\section{Двухсостоянийный носитель}

Чтобы два incidence-назначения стали состояниями одного объекта, введём
минимальный бинарный носитель $\mathcal H_{\rm bit}=\mathbb C^2$ с
\begin{equation}
 p_u=|u\rangle\langle u|,\qquad
 p_d=|d\rangle\langle d|,\qquad
 p_u+p_d=I_2,\qquad \sigma_z=p_u-p_d.
 \label{eq:v8-up-down-architecture-bit-projectors}
\end{equation}
Одномерный носитель не допускает двух ненулевых ортогональных проекторов,
тогда как $\mathbb C^2$ допускает их точно:
\begin{equation}
 \dim\mathcal H_{\rm bit}\ge2,\qquad
 \sigma_z^2=I_2.
 \label{eq:v8-up-down-architecture-minimality}
\end{equation}

Пусть $P_u,P_d$ --- Real-проекторы ранга шесть на двух цветных endpoint-
типах. Единый управляемый incidence-проектор равен
\begin{equation}
 \Pi_{\rm inc}=p_u\otimes P_u+p_d\otimes P_d,\qquad
 \Pi_{\rm inc}^2=\Pi_{\rm inc},\qquad \rank\Pi_{\rm inc}=12.
 \label{eq:v8-up-down-architecture-controlled-incidence}
\end{equation}
Бинарный носитель является gauge-синглетом и Real-чётным, поэтому
\begin{equation}
 [\Pi_{\rm inc},I_2\otimes\rho(g)]=0,\qquad
 J\Pi_{\rm inc}J^{-1}=\Pi_{\rm inc},\qquad
 [\Gamma,\Pi_{\rm inc}]=0.
 \label{eq:v8-up-down-architecture-compatibility}
\end{equation}
Так обе прежние диаграммы вложены в один типизированный операторный носитель,
а не сравниваются извне.

\section{Статический typed-гамильтониан}

Минимальный гамильтониан, наследующий квадратичный гиперзарядный индекс,
имеет вид
\begin{equation}
 H_{\rm inc}=\frac{25}{3}p_u+\frac{4}{3}p_d.
 \label{eq:v8-up-down-architecture-hamiltonian}
\end{equation}
После удаления общего сдвига он содержит единственный бинарный оператор:
\begin{equation}
 H_{\rm inc}=\frac{29}{6}I_2+\frac72\sigma_z,\qquad
 H_{\rm inc}^{0}=\frac72\sigma_z.
 \label{eq:v8-up-down-architecture-centered-hamiltonian}
\end{equation}
Для коэффициента $\gamma>0$ основное состояние единственно:
\begin{equation}
 E_u-E_d=7\gamma,\qquad
 \ker(\gamma H_{\rm inc}-E_{\min}I)=\mathbb C|d\rangle.
 \label{eq:v8-up-down-architecture-ground-state}
\end{equation}
Общий масштаб $\gamma$ не меняет выбранную ветвь, но его знак меняет её.
Поэтому условие $\gamma>0$ является частью условной архитектуры, пока
родительское происхождение знака не доказано.

При конечной обратной температуре $\beta$ Gibbs-состояние не является
чистым селектором:
\begin{equation}
 \frac{p_u^{(\beta)}}{p_d^{(\beta)}}=e^{-7\beta\gamma},\qquad
 p_u^{(\beta)}>0,\quad p_d^{(\beta)}>0,\qquad
 \lim_{\beta\gamma\to\infty}(p_u^{(\beta)},p_d^{(\beta)})=(0,1).
 \label{eq:v8-up-down-architecture-gibbs-limit}
\end{equation}

\section{Гамильтониан не создаёт релаксацию}

Чисто унитарная динамика
\begin{equation}
 \mathcal L_H(X)=i[\gamma H_{\rm inc},X],\qquad
 \operatorname{Fix}(\mathcal L_H)
 =\operatorname{span}\{p_u,p_d\},\qquad
 \dim\operatorname{Fix}(\mathcal L_H)=2
 \label{eq:v8-up-down-architecture-unitary-fixed-algebra}
\end{equation}
сохраняет обе популяции. Следовательно, уникальное основное состояние ещё
не означает динамический выбор.

Минимальный нуль-температурный dissipative completion использует один
скачок
\begin{equation}
 L_\downarrow=|d\rangle\langle u|,\qquad
 \mathcal L_\downarrow(X)=
 L_\downarrow^\dagger XL_\downarrow
 -\frac12\{L_\downarrow^\dagger L_\downarrow,X\}.
 \label{eq:v8-up-down-architecture-downward-jump}
\end{equation}
В картине состояний он действует на диагональном базисе как
\begin{equation}
 \mathcal L_{\downarrow *}(p_u)=p_d-p_u,\qquad
 \mathcal L_{\downarrow *}(p_d)=0.
 \label{eq:v8-up-down-architecture-population-flow}
\end{equation}
Его Heisenberg-супероператор имеет ранг три и одномерную неподвижную
алгебру:
\begin{equation}
 \rank\mathcal L_\downarrow=3,\qquad
 \ker\mathcal L_\downarrow=\mathbb CI_2,\qquad
 \rho_*=p_d.
 \label{eq:v8-up-down-architecture-primitive-relaxation}
\end{equation}
При конечной температуре необходимо добавить обратный скачок с отношением
скоростей
\begin{equation}
 \frac{\kappa_\uparrow}{\kappa_\downarrow}
 =e^{-7\beta\gamma}.
 \label{eq:v8-up-down-architecture-detailed-balance}
\end{equation}
Тогда стационарное состояние смешано; чистая $d$-ветвь получается только в
нулевом температурном пределе.

\section{Что построено и что не выведено}

Минимальные данные архитектуры разделяются на статические и динамические:
\begin{equation}
 \mathcal D_{\rm static}=\{\mathbb C^2,\Pi_{\rm inc},\gamma>0\},\qquad
 \mathcal D_{\rm dyn}=\{L_\downarrow,\kappa_\downarrow,
 \beta\ \text{или нуль-температурная граница}\}.
 \label{eq:v8-up-down-architecture-data-split}
\end{equation}
Ни бинарный носитель, ни управляемое endpoint-вложение, ни положительный
typed-коэффициент, ни скачок с физической скоростью не содержатся в старом
parent. Структура является минимальной достаточной архитектурой, но не
унаследованным выводом.

Итоговые реестры:
\begin{equation}
 N_{\rm typed\ static}=8/8,\qquad
 N_{\rm conditional\ relaxation}=5/5,\qquad
 N_{\rm inherited\ origin}=0/5.
 \label{eq:v8-up-down-architecture-ledgers}
\end{equation}

\begin{theorem}[Минимальный условный динамический селектор]
Двухсостоянийный gauge-тривиальный носитель является минимальным faithful-
представлением бинарного incidence-меню. Управляемый проектор
$\Pi_{\rm inc}$ совместим с gauge-, Real- и grading-структурами.
Гамильтониан квадрата гиперзаряда при $\gamma>0$ имеет единственное основное
состояние $|d\rangle$, но унитарная динамика популяции не выбирает. Один
скачок $|d\rangle\langle u|$ условно даёт уникальную стационарную $d$-ветвь.
\end{theorem}

\begin{nogocriterion}[Статический минимум не равен динамическому выбору]
Нельзя считать ветвь выбранной только из единственности основного состояния:
коммутатор с $H_{\rm inc}$ сохраняет обе диагональные популяции. Нельзя также
объявлять amplitude-damping скачок унаследованным, пока не построены его
системно-средовой оператор, состояние среды и физическая скорость.
\end{nogocriterion}

\begin{protocol}\begin{enumerate}
\item минимальная размерность бинарного носителя: \textbf{$2$};
\item управляемый incidence-проектор: \textbf{построен};
\item его ранг: \textbf{$12$};
\item gauge/Real/grading compatibility: \textbf{да};
\item typed-энергии: \textbf{$25/3$ и $4/3$};
\item статическая щель: \textbf{$7\gamma$};
\item основное состояние при $\gamma>0$: \textbf{$d$};
\item унитарная динамика выбирает популяцию: \textbf{нет};
\item минимальный downward jump: \textbf{$|d\rangle\langle u|$};
\item его неподвижная алгебра: \textbf{$\mathbb CI_2$};
\item скорость релаксации выведена: \textbf{нет};
\item typed static architecture: \textbf{$8/8$};
\item conditional relaxation: \textbf{$5/5$};
\item inherited origin: \textbf{$0/5$};
\item следующий гейт: \textbf{parent-origin бинарного носителя}.
\end{enumerate}\end{protocol}

Машинный сертификат сохранён в
\path{s2t_v8_baryon_c0_singlet_triplet_central_gap_isotypic_channel_extra_edge_mass_two_source_cubic_incidence_up_down_binary_selector_minimal_typed_discriminator_architecture_gate_results.json}.

\begin{thebibliography}{9}
\bibitem{v8-up-down-architecture-gks} V.~Gorini, A.~Kossakowski,
E.~C.~G.~Sudarshan, \emph{Completely Positive Dynamical Semigroups of
N-Level Systems}, J. Math. Phys. 17 (1976) 821.
\bibitem{v8-up-down-architecture-lindblad} G.~Lindblad,
\emph{On the Generators of Quantum Dynamical Semigroups},
Commun. Math. Phys. 48 (1976) 119.
\end{thebibliography}